% part: intuitionistic-logic
% chapter: propositions-as-types
% section: type-preservation

\documentclass[../../../include/open-logic-section]{subfiles}

\begin{document}

\olfileid{int}{pty}{tpr}

\olsection{પ્રકારનું સંરક્ષણ}

વિચારેલા ઘટાડા અને ક્રમચય-રૂપાંતરો
\Log{tN2i} તથા \Log{tN3i} તંત્રો માટે સીધા પણ વ્યાખ્યાયિત કરી શકાય.
\Log{tN3i} તંત્ર સંદર્ભ~$\Gamma$ની સાપેક્ષે
$\lambd$-પદ~$N$નો પ્રકાર એવું !!{formula}~$!A$ ગણે છે,
જે માટે $\Gamma \Sequent N : !A$ની
!!a{proof} \Log{tN3i}માં હોય.
$\lambd$-પદોના $\beta$-ઘટાડાના નિયમો
!!{proof}sના ઘટાડા અને ક્રમચય-રૂપાંતરોને બરાબર અનુરૂપ છે.
આ હકીકતનું એક મહત્ત્વનું પરિણામ છે:
સંદર્ભ~$\Gamma$ની સાપેક્ષે $\lambd$-પદનો પ્રકાર
$\beta$-રૂપાંતર પછી એ જ રહે છે.
આ ગુણધર્મને \emph{પ્રકારનું સંરક્ષણ} કહે છે.

\begin{prop}
$N$નો પ્રકાર $\Gamma$ની સાપેક્ષે~$!A$ હોય
અને $N \redone N'$ હોય,
તો $N'$નો પ્રકાર~$\Gamma$ની સાપેક્ષે~$!A$ છે.
\end{prop}

\begin{proof}
  $N$ અને $\Gamma \Sequent N : !A$ની
  !!{proof}s~$\delta$ની ઊંચાઈ પર આગમનથી
  આ સરળતાથી સ્થાપી શકાય:
  $M$ એ~$N$નું ઉપપદ હોય, તો $\delta$માં
  $\Gamma' \Sequent M : !B$થી પૂરી થતી
  ઉપ-!!{proof}~$\delta_1$ છે.
  $M$ સંકોચ્ય પદ હોય, તો~$\delta_1$ પર રૂપાંતર લાગુ પડે છે.
  $\delta_1$ પર ઘટાડો લગાવીને
  $\Gamma' \Sequent M' : !B$ની !!a{proof}~$\delta_1'$ મળે છે.
  \Log{tN3i}નાં રૂપાંતરો તપાસવાથી
  $M'$ એ~$M$ના સંકોચનનું પરિણામ છે તે દેખાય છે.
  $\delta$માં $\delta_1$ને બદલે $\delta_1'$ મૂકવાથી
  $\Gamma \Sequent N' : !A$ની !!a{proof}~$\delta'$ મળે છે;
  અહીં $N$માં ઉપપદ~$M$ને બદલે~$M'$ મૂકતાં $N'$ મળે છે.

  દાખલા તરીકે~\Intro{\land} પછી
  \Elim{\land} આવે તે માટેનું ઘટાડા-રૂપાંતર વિચારીએ:
  \[
  \bottomAlignProof
  \AxiomC{}
  \Deduce$\Gamma \fCenter N_1 : !B_1$
  \AxiomC{}
  \Deduce$\Gamma \fCenter N_2 : !B_2$
  \RightLabel{\Intro{\land}}
  \BinaryInf$\Gamma \fCenter \pair{N_1}{N_2} : !B_1 \land !B_2$
  \RightLabel{\Elim{\land}[i]}
  \UnaryInf$\Gamma \fCenter \proj{i}{\pair{N_1}{N_2}} : !B_i$
  \DisplayProof
  \quad\redone\quad
  \bottomAlignProof
  \AxiomC{}
  \Deduce$\Gamma \fCenter N_i : !B_i$
  \DisplayProof
  \]
  જમણી તરફની !!{proof}, એટલે $\delta_1'$નું સાબિતી-પદ
  $N_i$ છે તે દેખાય છે.
  તે બરાબર ડાબી તરફની !!{proof}~($\delta_1$)ના
  સાબિતી-પદ $\proj{i}{\pair{N_1}{N_2}}$ પર
  $\beta$-રૂપાંતર લગાવવાનું પરિણામ છે.

  અથવા \Intro{\lif} પછી \Elim{\lif} અનુમાન આવે
  અને તેના પર ઘટાડા-રૂપાંતર લગાવીએ તે વિચારીએ:
  \begin{multline*}
  \bottomAlignProof
  \AxiomC{$\delta_2$}
  \Deduce$x: !B_1, \Gamma \fCenter N : !B_2$
  \RightLabel{\Intro{\lif}}
  \UnaryInf$\Gamma \fCenter \lambd[\typeof{x}{!B_1}][N] : !B_1 \lif !B_2$
  \AxiomC{}
  \RightLabel{$\delta_3$}
  \Deduce$\Gamma \fCenter M : !B_1$
  \RightLabel{\Elim{\lif}}
  \BinaryInf$\Gamma \fCenter (\lambd[\typeof{x}{!B_1}][N] M) : !B_2$
  \DisplayProof
  \quad\redone\\
  \bottomAlignProof
  \AxiomC{}
  \RightLabel{$\delta_3$}
  \Deduce$\Gamma \fCenter M : !B_1$
  \RightLabel{$\Subst{\delta_2}{\delta_3}{x:!B_1}$}
  \Deduce$\Gamma \fCenter \Subst{N}{M}{x} : !B_2$
  \DisplayProof
  \end{multline*}
  અહીં પણ જમણી તરફની !!{proof}નું સાબિતી-પદ
  બરાબર ડાબી તરફની !!{proof}ના સાબિતી-પદ પર
  $\beta$-ઘટાડો લગાવવાનું પરિણામ છે.
  અલબત્ત,~$\delta_2$ની ઊંચાઈ પર આગમનથી
  સાવધાનીપૂર્વક આ તપાસવું પડે:
  $\delta_2$માં $\Gamma' \Sequent x:!B_1$ સ્વરૂપના બધા
  સ્વયંસિદ્ધોને બદલે $\Gamma \Sequent M:!B_1$ની
  !!{proof}~$\delta_3$ મૂકતાં ખરેખર
  $\Gamma \Sequent \Subst{N}{M}{x} : !B_2$ની
  !!a{proof} $\Subst{\delta_2}{\delta_3}{x:!B_1}$ મળે છે.
\end{proof}

!!{proof}s અને $\lambd$-પદો વચ્ચેની નિકટ અનુરૂપતાનું
બીજું પણ પરિણામ છે.
$N$ સંદર્ભ~$\Gamma$ની સાપેક્ષે પ્રકાર $!A$નું
$\lambd$-પદ હોય અને તેમાં સંકોચ્ય પદ~$M$ હોય
(એટલે તે સામાન્ય સ્વરૂપ ન હોય), તો તેને અનુરૂપ
$\Gamma \Sequent N : !A$ની \Log{tN3i} સાબિતી~$\delta$માં
$\Gamma' \Sequent M: !B$થી પૂરી થતી એવી
ઉપ-!!{proof} હોય છે જેના પર રૂપાંતર લગાવી શકાય.
$M$ને બદલે તેના સંકોચનનું પરિણામ મૂકવાથી
$N \redone N'$ હોય, તો~$\delta$ પર
અનુરૂપ ઘટાડો લગાવવાનું પરિણામ
$\Gamma \Sequent N' : !A$ની !!a{proof} છે.
એટલે સામાન્ય સ્વરૂપમાં ન હોય એવું દરેક $\lambd$-પદ
$\beta$-ઘટાડાથી બીજા પદમાં જાય છે.
આ ગુણધર્મને \emph{પ્રગતિ} કહે છે.

પ્રગતિ અને પ્રકારનું સંરક્ષણ મળીને ખાતરી આપે છે
કે $\lambd$-પદોનો $\beta$-ઘટાડો ક્યારેય ``અટકી જતો'' નથી:
પદ યોગ્ય રીતે પ્રકારિત હોય અને સામાન્ય સ્વરૂપમાં ન હોય,
તો તે બીજા યોગ્ય રીતે પ્રકારિત પદમાં ઘટે છે
(અને એ પદનો પ્રકાર એ જ હોય છે).

\end{document}

